% Issue 12 draft, 24 November 2026. All three Introduction paragraphs approved.
% Recruitment paragraph uses the previously agreed CIES wording.
% Final question/design/contribution paragraph: approved by the user.
% Total Introduction allowance: 0.40 pages; compiled fit remains for Issues 13--14.
\section*{Introduction}

Previous research has documented age differences in football recruitment. A 2022 CIES report found that 10.6\% of externally recruited players in Premier League squads had joined their club at age 30 or older \cite{poli2022recruitment}. Player performance typically peaks in the mid-to-late twenties \cite{dendir2016peak}, so a player of 30 or older is usually recruited after that peak. In this study, we compare Premier League joining rates among eligible recorded transfers involving players aged 30+ and those under 30.

Previous research has applied sentiment analysis to football coverage from multiple newspaper sources, including the Guardian \cite{dotgrau2019representation}. Large text corpora have also been used to measure how sports commentary describes different groups of players: Merullo et al.\ linked commentary to the players it named across 
American football broadcasts and found systematic differences by race \cite{merullo2019commentator}. We follow a similar player-linked design, applied to age. Monsees (2024) found that German and Swedish newspapers constructed football talent around similar themes, including age, but gave these themes different meanings across cultural contexts \cite{monsees2024talent}. These findings motivate investigating football transfers as a network combined with a linguistic perspective.

We ask: \textit{How does entry into the Premier League differ by player age, and how do the volume and tone of Guardian coverage around Premier League signings vary by age?} First, we compare the share of eligible transfers that end in the Premier League for players aged 30 or older and those under 30.
For the text analysis we link the same transfers of Premier League players to articles describing them. We use player names to find articles and select relevant sentences, and signing dates to define the 12-month windows before and after each move. We then compare how much coverage each age group receives. For signings with enough qualifying text in both periods, we describe the vocabulary of the selected sentences and test whether the change in tone over time differs between the same two age groups.





%The entry analysis asks where eligible transfers ended, so it includes moves to other leagues. The coverage analysis links selected Premier League signings to Guardian articles and compares tone before and after the move. We use the same age threshold in both samples and distinguish how much coverage a signing received from how its tone changed. %The selected sentences from these articles form the corpus for the sentiment analysis. To describe its vocabulary, we tokenize the sentences and examine which tokens occur most often and how their frequencies vary by rank.
